\section{Additional continuous geometry proofs}
\label{app:geometry}

\subsection{Uniform coordinate-sheet bounds}

\begin{lemma}[Graph geometry and coordinate multiplicity]
\label{app:geometry:lem-graphs}
Let $S=\{x+\phi(x):x\in U\}$, where $U$ is a convex subset of a
$d$-dimensional linear space $T$ and $\phi:U\to T^\perp$ is $C^2$ with
$\|D^2\phi\|\leq H$. Then
\begin{equation}\label{app:geometry:eq-two-point}
 \dist(z-y,T_yS)\leq\frac H2|z-y|^2\qquad(y,z\in S).
\end{equation}
More generally, if a manifold has this inequality with $H=1/\rho$, then
each dominant coordinate projection in the definition preceding
\cref{geom:lem-arcsine} is injective on subsets of diameter less than
\[
 r_\rho=\frac{2\rho}{1+\binom nd^{1/2}}.
\]
In particular its $M(S)$ is at most $\binom nd$ when its diameter is
less than $r_\rho$. Uniformly bounded graph Hessians give a uniform
choice of this diameter. A zero Hessian permits any diameter.
\end{lemma}
\begin{proof}
For $y=x+\phi(x)$ and $z=x'+\phi(x')$, the vector
$(x'-x)+D\phi(x)(x'-x)$ lies in $T_yS$. Taylor's theorem on the segment
in $U$ bounds the remainder by $H|x'-x|^2/2\leq H|z-y|^2/2$.

Suppose $y,z\in S^I$ have $y_I=z_I$. Decompose $z-y=t+\nu$ into its
tangent and normal parts at $y$. The smallest singular value of
$\pi_I|_{T_yS}$ is at least $J_I(y)$: all singular values are at most
one and their product is $J_I$. Consequently
\[
 |t|\leq\binom nd^{1/2}|\pi_I t|
       =\binom nd^{1/2}|\pi_I\nu|
       \leq\binom nd^{1/2}|\nu|.
\]
The two-point inequality gives
$|z-y|\leq(1+\binom nd^{1/2})|z-y|^2/(2\rho)$.
Either $y=z$ or their distance is at least $r_\rho$. This proves the
multiplicity bound.
\end{proof}

\subsection{Proof of the nonlinear constraint transfer}
\label{app:geometry:proof-tube-kkt}

We prove \cref{geom:thm-tube-kkt} with the exact/relaxed constraint tuple
specified there. Root bounds are exactly kept inequalities. Let $E$
collect all exact equalities and active exact inequalities, including
active root coordinates. LICQ supplies a vector $e$ whose derivatives
are zero along these constraints and every other active constraint
except one selected relaxed constraint. Set its derivative to one for
a relaxed inequality with positive multiplier, or to the sign of the
multiplier for a relaxed equality with nonzero multiplier. Normalize to
a unit vector $e_0$. KKT stationarity gives
\[
 E'(z_*)e_0=0,\qquad \nabla f(z_*)\cdot e_0=-\sigma<0.
\]
Full row rank gives a fixed right inverse $R$ of $E'(z_*)$. When $E$ is
empty no correction is needed. Otherwise apply the $C^2$ implicit
function theorem to
\[
 H(y,t,\xi)=E(y+t e_0+R\xi)-E(y).
\]
Its derivative in $\xi$ at $(z_*,0,0)$ is the identity. On an open
product neighborhood there is a $C^2$ solution $\xi(y,t)$, with
$\xi(y,0)=0$ by uniqueness. Define
\[
 Z(y,t)=y+t e_0+R\xi(y,t).
\]
Then $E(Z(y,t))=E(y)$ and $Z(y,0)=y$. Differentiating at $(z_*,0)$
gives $\xi_t(z_*,0)=0$, and hence $Z_t(z_*,0)=e_0$.
Also $Z_y(y,0)=I$ throughout a smaller neighborhood.

On a product neighborhood with compact closure, shrink the spatial and
time ranges to obtain constants $\mu,c,t_0>0$, $C\geq0$ such that
for feasible $y$ and $0\leq t\leq t_0$,
\begin{align}
 |Z_t(y,t)|&\leq2,&
 f(Z(y,t))&\leq f(y)-\mu t+Ct^2,\label{app:geometry:eq-descent}\\
 |Z(y,t)-y|&\leq2t,&v(Z(y,t))&\leq ct.
 \label{app:geometry:eq-violation}
\end{align}
Continuity of the strictly negative time derivative at $(z_*,0)$ and
a bound on the second derivative of $f\circ Z$ give the descent
estimate. Bounded derivatives of the original relaxed functions and
their nonpositive or zero values at $y$ give the violation estimate.
Active exact inequalities preserve their feasible value at $y$, exact
equalities preserve zero, and active root bounds preserve their
coordinate value. All inactive exact restrictions have uniform slack
on the smaller neighborhood and retain it for small $t$. Thus
$Z(y,t)\in D\cap P_{\mathrm{ex}}$.

For the covering conclusion, put $a=\eta+\epsilon$ and
$t=3a/\mu$. If $a$ is small enough that $t\leq t_0$ and $Ct^2\leq a$,
then for $y\in E(\eta)$,
\begin{equation}\label{app:geometry:eq-superoptimal}
 f(Z(y,t))\leq f_*+\eta-3a+a
                   =f_*-\epsilon-a<f_*-\epsilon.
\end{equation}
The tube dichotomy therefore forces
\[
 \beta q_{C_e}(Z(y,t))<v(Z(y,t))\leq Aa,\qquad A=3c/\mu
\]
in every owner containing the shifted point. It lies within
$r=\sqrt{Aa/\beta}$ in sup norm of a test vertex. If also
$a\leq c\mu/(12\beta)$, its displacement $6a/\mu$ is at most $r$.
The original $Y$ is consequently covered by at most $2^nK$ cubes of
side $4r$. This proves \cref{geom:eq-local-tube-cover}, and one may take
the effective coefficient $\alpha_{\mathrm{eff}}=\mu\beta/(12c)$ in
the side length $2\sqrt{a/\alpha_{\mathrm{eff}}}$.

For the integral conclusion, LICQ gives a feasible active-stratum graph
\[
 y(x)=z_*+x+\phi(x),\qquad x\in B_r\subset T,
 \qquad \phi(0)=D\phi(0)=0,
\]
where $T$ is its $d$-dimensional tangent space. The graph is feasible:
all active inequalities are zero and inactive ones retain their slack.
Put $m_T(x)=m(y(x))$. Because $z_*$ minimizes on this graph,
$m_T(0)=Dm_T(0)=0$. Define
\[
 W_\epsilon(x)=Z\left(y(x),\frac{3(m_T(x)+\epsilon)}\mu\right),
 \qquad
 \Psi_\epsilon(x)=\pi_T(W_\epsilon(x)-z_*).
\]
These functions are jointly $C^2$ in an open neighborhood of $(0,0)$;
the time range can be kept inside the IFT domain. At $(0,0)$ the
derivative of the time parameter vanishes, so
$D_x\Psi_0(0)=I$. Choose a fixed ball $B_r$ and a fixed compact
parameter range $0\leq\epsilon\leq\epsilon_0$ such that
\begin{equation}\label{app:geometry:eq-shift-inverse}
 \|D_x\Psi_\epsilon-I\|\leq\tfrac12,
                 \qquad |\Psi_\epsilon(0)|\leq r/8.
\end{equation}
The mean-value formula along line segments gives
$|\Psi_\epsilon(x)-\Psi_\epsilon(x')|\geq|x-x'|/2$.
Thus $\Psi_\epsilon$ is injective and its determinant is at least
$2^{-d}$ in absolute value.

Every image contains the common ball $V=B_{r/4}\subset T$. Indeed, for
$u\in V$, the map $x\mapsto u-(\Psi_\epsilon(x)-x)$ is a contraction
of the closed ball $B_r$ into itself: its norm is at most
$r/4+r/8+r/2<r$. Its fixed point solves $\Psi_\epsilon(x)=u$.
The inverse on $V$ has derivative norm at most two. Differentiating
the inverse identity gives
\[
 D^2\Psi_\epsilon^{-1}
 =-(D\Psi_\epsilon)^{-1}D^2\Psi_\epsilon
      [D\Psi_\epsilon^{-1},D\Psi_\epsilon^{-1}],
\]
so its second derivatives are bounded uniformly in $\epsilon$.
The shifted image over $V$ is therefore a graph with uniformly bounded
second derivatives:
\[
 u\longmapsto z_*+u+
       \pi_{T^\perp}(W_\epsilon(\Psi_\epsilon^{-1}(u))-z_*).
\]
Restrict the original graph to a still smaller fixed patch $S$ whose
shifted projections lie in $V$. Its shifted images have uniformly
small diameter, by the bounded derivative of $W_\epsilon$ and small
$\epsilon_0$. \Cref{app:geometry:lem-graphs} now gives
$M(Z_\epsilon(S))\leq\binom nd$ for the maps
$Z_\epsilon(y)=Z(y,3(m(y)+\epsilon)/\mu)$.
This argument controls multiplicity uniformly even though the image
need not be a translate of the original patch.

The area factor of $W_\epsilon$ is at least that of its tangential
projection, hence at least $2^{-d}$. The input graph's area factor is
bounded above by $B_\phi=(1+\sup\|D\phi\|^2)^{d/2}$.
Consequently the Jacobian relative to stratum measure satisfies
\[
 J_{Z_\epsilon}\geq j_*:=2^{-d}/B_\phi>0.
\]
Finally shrink $S$ and $\epsilon_0$ so that the descent remainder is
at most $m(y)+\epsilon$ throughout. The previous strict-descent
calculation gives
\[
 f(Z_\epsilon(y))<f_*-\epsilon,
              \qquad v(Z_\epsilon(y))\leq A(m(y)+\epsilon).
\]
For each owner, the dichotomy implies
$\beta q_{C_e}(Z_\epsilon(y))<A(m(y)+\epsilon)$ on its preimage.
Use the injective area formula and the lower Jacobian, followed by
\cref{geom:lem-arcsine} on the shifted patch, to get
\[
 \int_{\{y\in S:Z_\epsilon(y)\in A_e\}}
        (m(y)+\epsilon)^{-d/2}\,d\mathcal H^d(y)
 \leq j_*^{-1}(A/\beta)^{d/2}C_{n,d}\binom nd.
\]
Summing over a covering family proves the first inequality in
\cref{geom:eq-tube-integral}. For a run the all-owner family is the
one that covers these images; owners meeting $F$ alone need not cover
them.

Taylor's theorem gives $0\leq m_T(x)\leq H|x|^2$ on a fixed parameter
ball, with $H\geq1$. Its input graph area factor is at least one, so
\[
 \int_S(m+\epsilon)^{-d/2}\,d\mathcal H^d
 \geq d\omega_d\int_0^{r'}
          t^{d-1}(Ht^2+\epsilon)^{-d/2}\,dt,
\]
where $\omega_d$ is the unit-ball volume. For
$t\geq\sqrt{\epsilon/H}$ the integrand is at least
$(2H)^{-d/2}/t$. The resulting logarithm proves the last assertion.

\subsection{Why a single covering scale is insufficient}

\begin{example}[Many near-optimal wells]\label{app:geometry:ex-wells}
Let $D=F=[0,1]$, $t=y-1/2$, and, for an integer $K\geq2$, define
\[
 f_K(y)=K^{-2}(1-\cos(2\pi Kt))
                  +K^{-2}(1-e^{-K^2t^2}).
\]
Its optimum is zero at $t=0$. Use
$f_B=f_K-32q_B$ and $R_B=B$. This is a convex underestimator because
\[
 f_K''(y)=4\pi^2\cos(2\pi Kt)+(2-4K^2t^2)e^{-K^2t^2}
                         \geq-4\pi^2-4e^{-3/2}>-64.
\]
Thus the lower and vertex-vanishing upper coefficients are both $32$.
The first and second derivatives are uniformly bounded on $D$, so the
remaining Lipschitz and second-order constants can also be fixed across
the family; feasibility imposes no error-bound difficulty.

For $|t|\leq1/(2K)$, $1-\cos(2\pi Kt)\geq8K^2t^2$, so
$f_K\geq8t^2$. For $|t|\geq1/(2K)$ the second term is at least
$(1-e^{-1/4})K^{-2}>0.22K^{-2}$. Hence for
$\epsilon_K=0.2K^{-2}$,
\[
 E(\epsilon_K)\subset
           \{|t|\leq\sqrt{\epsilon_K/8}\}.
\]
For any fixed $c>0$, its cover by side-$c\sqrt{\epsilon_K}$ intervals
has size at most $2+1/(\sqrt2c)$, independently of $K$.

The points $t=k/K$, $|k|\leq K/2$, number at least $K$ and all have
$f_K\leq K^{-2}$. In \cref{geom:thm-cover-lower} take
$\eta=K^{-2}$. Its covering interval has side
$2\sqrt{1.2/32}/K<1/K$, so it covers at most one of these points.
Consequently $N_{\mathrm{cover}}(\epsilon_K)\geq K/2$.
The ratio cannot be bounded by a uniform constant times
$\log(1/\epsilon_K)$ times the single-scale covering count, since that
expression is $O(\log K)$.
\end{example}

\subsection{Proof of the regular KKT rates}
\label{app:geometry:proof-kkt-rates}

Fix a global minimizer $z_*$. In its local feasible description, write
$A$ for the active inequalities, $\mu_j>0$ for their multipliers,
and $\lambda_k$ for equality multipliers. Let
$\mathcal L=f+\sum_{j\in A}\mu_jg_j+\sum_k\lambda_k h_k$, with all
constraints zero at the minimizer. KKT gives $\nabla\mathcal L(z_*)=0$.
For feasible $y=z_*+\delta$, Taylor's theorem gives
\begin{equation}\label{app:geometry:eq-kkt-growth}
 m(y)=\tfrac12\delta^\top H_{\mathcal L}\delta+o(|\delta|^2)
                              +\sum_{j\in A}\mu_j(-g_j(y)).
\end{equation}
If local quadratic growth failed, there would be feasible $y_k\to z_*$
with $m(y_k)/|\delta_k|^2\to0$, and, after a subsequence,
$\delta_k/|\delta_k|\to u$ with $|u|=1$. Boundedness of the Hessian and
nonnegativity of the slacks in \cref{app:geometry:eq-kkt-growth} imply
$-g_j(y_k)=O(|\delta_k|^2)$ for each $j\in A$.
Taylor expansion of these constraints and the equalities gives
$u\in T$, the common nullspace of their gradients. If there are no
active inequalities, the equalities alone give this conclusion.
Positive definiteness of $H_{\mathcal L}$ on $T$ then implies
\[
 \liminf_k\frac{m(y_k)}{|\delta_k|^2}
            \geq\tfrac12 u^\top H_{\mathcal L}u>0,
\]
a contradiction. If $T=\{0\}$, already $u\in T$ contradicts $|u|=1$.
Thus each minimizer has local quadratic growth. Disjoint neighborhoods
of the finitely many minimizers, together with the positive minimum of
$m$ on their compact feasible complement, give global quadratic growth
to the optimal set. The upper part of \cref{geom:prop-dimension} gives
the logarithmic upper bound.

For a minimizer whose active stratum has dimension $d\geq1$, LICQ
gives an injective $C^2$ coordinate graph $y(x)$ of that stratum.
Inactive constraints have slack on a small patch, so the graph is
feasible. All active constraints vanish there and hence
$m(y(x))=\mathcal L(y(x))-f_*\leq C|x|^2$. Choose a $d$-coordinate
projection invertible at the minimizer and shrink to a graph patch
with bounded area factor. For every box, the arithmetic--geometric
mean and arcsine calculation in \cref{geom:lem-arcsine} bounds
$\int_{S\cap B}q_B^{-d/2}$ by a fixed constant using this projection
alone. Also the radial calculation at the end of
\cref{app:geometry:proof-tube-kkt} gives
$\int_S(m+\epsilon)^{-d/2}\geq c\log(1/\epsilon)$.
Validity therefore gives the objective-gap logarithmic lower bound.
The tube lower bound is \cref{geom:thm-tube-kkt} for its compatible
original tuple.

For a unique zero-dimensional active stratum, the active gradient
matrix is square and nonsingular. Consequently
\[
 |\delta|\leq C\left(\sum_{j\in A}|\nabla g_j(z_*)\delta|
                  +\sum_k|\nabla h_k(z_*)\delta|\right).
\]
Taylor expansion and feasibility show
\[
 \sum_{j\in A}(-g_j(y))\geq c|\delta|-C'|\delta|^2.
\]
If $A$ is nonempty, positive multipliers and
\cref{app:geometry:eq-kkt-growth} give local sharp growth
$m(y)\geq c_s|y-z_*|$ after shrinking the neighborhood. If $A$ is
empty, the equalities give $|\delta|\leq C'|\delta|^2$, so $z_*$ is
locally isolated in $F$; sharp growth holds vacuously there.
Uniqueness and compactness extend sharp growth globally in both cases.
This handles the empty-active-inequality case without an undefined
minimum of multipliers.

To build an exact certificate under \cref{geom:eq-upper-q}, choose a
coordinate grid of side at most $s$ through $z_*$, clipped at $D$.
Take $s$ small enough that $\alpha'ns^2/4\leq v_0$.
For $z\in R_C$, a nearest feasible $y$ and sharp growth give
\[
 f_C(z)-f_*
 \geq c_s|z-z_*|-K'q_C(z),\qquad
 K'=\alpha'(1+L_f\kappa+c_s\kappa).
\]
Every grid interval has $z_{*,i}$ outside it or at an endpoint, so
$d_i^C(z)\leq|z_i-z_{*,i}|$. Hence
\[
 q_C(z)\leq s\sum_i d_i^C(z)
                  \leq s\sqrt n\,|z-z_*|.
\]
Choosing also $s\leq c_s/(K'\sqrt n)$ gives $f_C(z)\geq f_*$ on
every relaxation. If $\alpha'=0$, the error bound gives feasible
relaxed points and the root is exact directly. The finite clipped
grid is a split-tree partition and proves the exact-certificate claim.

Finally, if a minimizer coordinate stays a relative distance at least
$\delta\in(0,1/2]$ from both endpoints of its dyadic interval, the
objective lower gap gives
$L(B)\leq f_*-\alpha\delta(1-\delta)s_j^2$ for a box containing it.
Such a box remains unpruned while this quantity exceeds $\epsilon$,
giving a logarithmic number of processed levels. For a coordinate
position $p/q$ in lowest terms with $q>1$ odd, every dyadic fractional
position is between $1/q$ and $1-1/q$. This establishes the stated
example of the middle-position condition.

\subsection{Convex centroids and fractional covers}

\begin{lemma}[Centroid distance from supporting planes]
\label{app:geometry:lem-centroid}
Let $K$ be a bounded convex subset of a $p$-flat with positive
relative volume. Its centroid belongs to $K$. For any linear
functional of width $W$ on $K$, its centroid is at distance at least
$W/(p+1)$ in functional value from each supporting plane.
The assertion holds when $K$ is partly open.
\end{lemma}
\begin{proof}
The closure has the same volume and first moments, since a convex
body's boundary has zero relative volume. Its centroid is in its
relative interior: a supporting hyperplane through a putative boundary
centroid would put a nonnegative, nonzero linear functional above its
minimum on a positive-volume subset, contradicting that centroid's
value. The relative interiors of a convex set and its closure coincide
and belong to the original convex set, so the centroid belongs to $K$.

Work with the closure and normalize the functional's minimum to zero.
If $W=0$ the conclusion is immediate. Choose an apex $a$ at its maximum
$W$. For $0\leq t\leq W$, the homothetic image
\[
 (1-t/W)\cl K+(t/W)a
               \subset\cl K\cap\{h\geq t\}
\]
has volume $(1-t/W)^p\vol_p(K)$. Integrating these tail volumes gives
the mean functional value at least
$\int_0^W(1-t/W)^p\,dt=W/(p+1)$.
Apply the same argument to $W-h$ for the other supporting plane.
\end{proof}

\begin{lemma}[Half-integral fractional vertex covers]
\label{app:geometry:lem-half-integral}
Every finite graph has an optimal fractional vertex cover with entries
in $\{0,1/2,1\}$.
\end{lemma}
\begin{proof}
Truncating an entry greater than one to one preserves every edge
constraint and improves the objective. Thus an optimum exists in the
compact polytope $0\leq z_i\leq1$, $z_i+z_j\geq1$. Choose an optimal
vertex of that polytope. On the vertices with $0<z_i<1$, form the graph
of tight equations $z_i+z_j=1$. A bipartite component admits a small
perturbation adding $t$ on one side and subtracting it on the other.
All its tight equations remain satisfied; every other constraint
involving these vertices is slack, so both signs of a sufficiently
small perturbation stay feasible. This contradicts extremality.
Every fractional component therefore has an odd cycle. Its tight
equations force all entries around the cycle, and hence throughout
the component, to be $1/2$. All remaining entries are zero or one.
\end{proof}

\subsection{Proof of the two-dimensional aligned certificate}
\label{app:geometry:proof-aligned-two}

We prove \cref{geom:prop-aligned-two}. On the optimal segment,
$q(y):=g(a,y)=f_*-cay$ is affine. For a feasible transverse sign
$\sigma\in\{-1,1\}$ define
\[
 h_\sigma(y)=\lim_{t\downarrow0}
             \frac{g(a+\sigma t,y)-q(y)}t,
                         \qquad L_y<y<U_y.
\]
Convexity in $x$ gives existence of the derivative. The minimum
condition for $m$ along the feasible transverse direction gives a
finite lower bound; a fixed positive secant gives a finite upper
bound. Thus these derivatives are finite on the open interval.

For two interior values $y_0,y_1$, choose $0<\lambda<1$ and interior
$y_2$ with $y_1=\lambda y_0+(1-\lambda)y_2$. For small $t>0$ convexity
gives
\[
 g(a+\sigma t,y_1)
 \leq\lambda g(a+\sigma t/\lambda,y_0)+(1-\lambda)q(y_2).
\]
Subtract the affine identity for $q(y_1)$, divide by $t$, and let
$t\downarrow0$. The right-hand difference quotient is exactly the
one defining $h_\sigma(y_0)$, so
$h_\sigma(y_1)\leq h_\sigma(y_0)$. Reversing the roles gives equality.
Each applicable transverse derivative of $g$ is therefore constant
along the open segment. Write the nonnegative transverse derivatives
of $m$ as
\[
 r_+(y)=h_++cy,\qquad r_-(y)=h_--cy.
\]
For $c>0$, their nonnegativity throughout the interval forces
\[
 r_+(y)\geq c(y-L_y),\qquad
 r_-(y)\geq c(U_y-y).
\]
Convexity of $x\mapsto m(x,y)$ with $m(a,y)=0$ then gives
\[
 m(x,y)\geq c(x-a)(y-L_y)\quad(x\geq a),\qquad
 m(x,y)\geq c(a-x)(U_y-y)\quad(x\leq a).
\]
On the respective child boxes these expressions are among the endpoint
products whose minimum is the positive-coefficient gap in
\cref{geom:eq-edge-gap}. Hence $\Gamma_B\leq m$ on both children.
For $c<0$, the same nonnegativity argument gives
\[
 r_+(y)\geq |c|(U_y-y),\qquad
 r_-(y)\geq |c|(y-L_y),
\]
which supplies the negative-coefficient gap products instead.
Continuity extends the inequalities to the segment endpoints and the
closed box boundaries. At an endpoint $a$ only the feasible sign is
needed and its box is the root. The pointwise pruning criterion
\cref{geom:eq-mcc-valid} proves exactness in every case.
